\documentclass[a4paper,12pt,fleqn]{article}%fleqn pour aligner les equations à gauche
\usepackage{../../Styles/Exercice}


\newcommand{\chnum}{Chapitre 3}
\newcommand{\chtitle}{Oxydo-réduction et Avancement}
\newcommand{\dtype}{Exercices}
  
  
\begin{document}
\maketitle
\section{Espèces conjuguées}
\begin{enumerate}
    \item \begin{enumerate}[label = \alph*]
        \item \underline{\ce{Br2(aq)}} \ce{+ 2 e- = } \underline{\ce{2 Br-(aq)}}
        \item \underline{\ce{Sn^2+(aq)}} \ce{+ 2 e- = }\underline{\ce{Sn(s)}}
        \item \underline{\ce{MnO2(s)}}\ce{+ 4 H+(aq) + 2 e- =} \underline{\ce{Mn^2+(aq)}} \ce{+ 2 H2O(l)}
    \end{enumerate}
    \item  \begin{enumerate}[label = \alph*]
        \item \ce{Br2}/\ce{Br-}
        \item \ce{Sn^2+}/\ce{Sn}
        \item \ce{MnO2}/\ce{Mn^2+}
    \end{enumerate}
\end{enumerate}

\section{Ecrire l'équation d'une réaction redox}
\begin{enumerate}[label = \alph*]
    \item \ce{Cu(s) + Br2(aq) -> Cu^2+(aq) + 2Br-(aq)}
    \item \ce{2I-(aq) + 2Fe^3+(aq) -> I2(s) + 2Fe^2+(aq)}
    \item \ce{2S2O3^2-(aq) + I2(s) -> S4O6^2-(aq) + 2I-(aq)}
    \item \ce{2H+(aq) + Fe(s) -> H_2(g) + Fe^2+(aq)}
    \item \ce{2Ag+(aq) + Cu(s) -> 2Ag(s) + Cu^2+(aq)}
    \item \ce{MnO4-(aq) + 5Fe^2+(aq) + 8H+(aq) -> Mn^2+(aq) + 5Fe^3(aq) + 4H2O(l)}
\end{enumerate}

\section{Identification d'un oxydant et d'un réducteur}
\begin{enumerate}
    \item \begin{enumerate}[label = \alph*]
        \item \fbox{\ce{2Ag+(aq)}} + \fcolorbox{blue}{white}{\ce{H2(g)}} \ce{->} \ce{2Ag(s) + 2H+(aq)}
        \item \fbox{\ce{S2O8^2-(aq)}} + \fcolorbox{blue}{white}{\ce{Cu(s)}} \ce{-> 2SO4^2-(aq) + Cu^2+(aq)}
        \item \fbox{\ce{Au^3+(aq)}} + \fcolorbox{blue}{white}{\ce{3Fe^3+(aq)}} \ce{-> Au(s) + 3Fe^3+(aq)}
    \end{enumerate}
    \item \begin{enumerate}[label = \alph*]
        \item \ce{Ag+(aq)}/\ce{Ag(s)}  \hspace{2em} \ce{H+(aq)}/\ce{H2(g)}
        \item \ce{S2O8^2-(aq)}/\ce{SO4^2-(aq)}  \hspace{2em} \ce{Cu^2+(aq)}/\ce{Cu(s)}
        \item \ce{Au^3+(aq)}/\ce{Au(s)} \hspace{2em} \ce{Fe^3+(aq)}/\ce{Fe^2+(aq)}
    \end{enumerate}
\end{enumerate}

\section{Réaction ou pas réaction ?}
\begin{enumerate}
    \item Mélange 1 : \ce{Au^3+(aq)} est oxydant, \ce{Fe(s)} est réducteur : \emph{Réaction}\\
    Mélange 2 : \ce{Cr^3+(aq)} est réducteur, \ce{Cr2O7^2-(aq)} est réducteur mais ils appartiennent au même couple : \emph{Pas de réaction}\\
    Mélange 3 :\ce{Ag+(aq)} est un oxydant, \ce{Fe^2+(aq)} est un réducteur : \emph{Réaction}\\
    Mélange 4 : \ce{Mn^2+(aq)} est un réducteur, \ce{Zn(s)} est un réducteur : \emph{Pas de réaction}
    \item Mélange 1 : \ce{2Au^3+(aq) + 3Fe(s) -> 2Au(s) + 3Fe^3+(aq)}\\
    Mélange 3 : \ce{Ag+(aq) + Fe+(aq) -> Ag(s) + Fe^3+(aq)}
\end{enumerate}

\section{L'ion nitrate dans un engrais}
\begin{enumerate}
    \item Couples mis en jeu : \\
    \ce{NO3-(aq)}/\ce{NO(g)} et \ce{Fe^3+(aq)}/\ce{Fe^3+(aq)}\\
    \ce{NO^3-(aq)} ne peut que jouer un rôle d'oxydant, il faudra donc choisir le couple où \ce{Fe^2+(aq)} joue le rôle de réducteur pour qu'il puisse y avoir une réaction chimique.
    \item Equation de la réaction $$\ce{NO3-(aq) + 4H+(aq) + 3e- = NO(g) + 2H2O(l)}$$
    $$ 3\times \left(\ce{Fe^2+(aq) = Fe^3+(aq) + e-}\right)$$
    \begin{center}
        \rule{0.5\textwidth}{.4pt}
    \end{center}
    $$ \ce{NO3-(aq) + 4H+(aq) + 3Fe^2+(aq) -> NO(g) + 2H2O(l) + 3Fe^3+(aq)}$$
\end{enumerate}

\section{Trouver une équation d'oxydoréduction}
On observe, au début de la réaction, que la solution est violette, ce qui indique la présence d'ions \ce{MnO4- (aq) dans la solution}. Les photos montrent une solution qui se décolore mais reste limpide (absence de trouble ou de précipité). Cela indique que le produit de na réaction ne peut pas être \ce{MnO2(s)} mais plutôt \ce{Mn^2+ (aq)}. L'équation proposée n'est pas la bonne. Il faudrait plutôt écrire :
$$2\times \left(\ce{MnO4- (aq) + 8H+ (aq)+ 5e- = Mn^2+(aq) + 4H2O(l)}\right)$$
$$5\times \left(\ce{H2C2O4(aq) = 2CO2(g) + 2 H+(aq) + 2 e-}\right) $$
\begin{center}
    \rule{0.5\textwidth}{.4pt}
\end{center}
$$ \ce{2MnO4- (aq) + 16H+(aq) + 10e- + 5H2C2O4(aq) -> 2Mn^2+(aq) + 8H2O(l) + 10CO2(g) + 10H+(aq) + 10e-}$$
$$ \ce{2MnO4- (aq) + 6H+(aq) + 5H2C2O4(aq) -> 2Mn^2+(aq) + 8H2O(l) + 10CO2(g)}$$

\section{Du plomb dans l'eau}
Pour que le plomb se retrouve dans l'eau, il faut qu'il puisse réagit avec du dioxygène dissout dans l'eau. Le propriétaire décide de refaire ses canalisations et dispose de deux choix : le cuivre ou le zinc.\\
Dans le cas du cuivre, on remarque que le cuivre solide (\ce{Cu(s)}) ont un poitntiel $E^0$ plus élevé ce qui en fait un mauvais réducteur vis à vis des ions \ce{Pb^2+(aq)} potentiellement présent dans le réseau de la ville.\\
Si il choisit le zinc, le potentiel du zinc solide (\ce{Zn(s)}) est inférieur à celui du couple du plomb ce qui favorisera la réaction avec les ions \ce{Pb^2+(aq)}.\\
La rénovation doit prendre en compte les propriétés d'oxydoréduction du matériau utilisé pour la tuyauterie : le cuivre serait dangereux à utiliser alors que le zinc permettrait de limiter la quantité d'ions \ce{Pb^2+(aq)} dans l'eau.\\
On notera qu'il devrait faire une demande aux services municipaux car les canalisations en plomb sont interdites depuis les années 1960 et qu'elles ne sont donc pas aux normes.

\section{Bouteille bleue}
\begin{enumerate}
    \item Demi-équations : $$\ce{C6H11O7-(aq) + 3H+(aq)  2e- = C6H12O6(aq) + H2O(l)}$$
    $$ \ce{O2(g) + 2H2O(l) + 4e- = 4HO-(aq) }$$
    $$ \ce{B^2+(aq) + H+(aq) + 2e- = BH+(aq)}$$
    \item Les ions \ce{B2+(aq)} sont à l'origine de la couleur bleue de la solution. Leur disparition vient de la réduction de ces ions en milieu basique par \ce{HO-(aq)} pour former \ce{BH+(aq)} et \ce{O2(g)}.
    \item Lors de l'agitation on dissout du dioxygène dans la solution ce qui permet la réaction avec le glucose : $$ \ce{2C6H12O6(aq) + O2(g) + 4H2O(l) -> 2C6H11O7-(aq) + 6H+(aq) + 4HO-(aq)}$$
    Les ions \ce{C6H11O7-(aq)} formés peuvent alors réagir avec \ce{BH+(aq)} pour former des ions \ce{B^2+(aq)} bleus : $$ \ce{C6H11O7-(aq) + 3H+(aq) + BH+(aq) -> C6H12O6(aq) + H2O(l) +B^2+(aq) + H+(aq)}$$
    \item En théorie, cette réaction pourrait être réalisée un nombre indéfini de fois. Toute fois des paramètres comme la solubilité du dioxygène dans l'eau, le pH de la solution et les constantes d'équilibre limitent le nombre de fois où elle peut être répétée.
\end{enumerate}
\end{document}
